\documentclass[12pt, titlepage]{article}

\usepackage{booktabs}
\usepackage{tabularx}
\usepackage{hyperref}
\hypersetup{
    colorlinks,
    citecolor=black,
    filecolor=black,
    linkcolor=red,
    urlcolor=blue
}
\usepackage[round]{natbib}

\input{../Comments.text}
\input{../Common.text}

\begin{document}

\title{Verification and Validation Report: \progname}
\author{\authname}
\date{\today}

\maketitle

\pagenumbering{roman}

\section{Revision History}

\begin{tabularx}{\textwidth}{p{3cm}p{2cm}X}
\toprule {\bf Date} & {\bf Version} & {\bf Notes}\\
\midrule
April 22, 2026 & 1.0 & Initial version \\
\bottomrule
\end{tabularx}

~\newpage

\section{Symbols, Abbreviations and Acronyms}

\renewcommand{\arraystretch}{1.2}
\begin{tabular}{l l}
  \toprule
  \textbf{symbol} & \textbf{description}\\
  \midrule
  T & Test\\
  FR & Functional Requirement\\
  NFR & Nonfunctional Requirement\\
  SRS & Software Requirements Specification\\
  MG & Module Guide\\
  MIS & Module Interface Specification\\
  VnV & Verification and Validation\\
  CI & Continuous Integration\\
  MNE & MNE-Python library for M/EEG processing\\
  \bottomrule
\end{tabular}\\

\newpage

\tableofcontents

\listoftables

\listoffigures

\newpage

\pagenumbering{arabic}

This document reports the verification and validation outcomes for \progname{}.
The report is based on the approved VnV strategy and test definitions in the VnV Plan, and traces results back to the requirements (SRS) and design artifacts (MG and MIS).

\section{Functional Requirements Evaluation}

Functional requirements in the SRS are evaluated using system tests FR-01 to FR-04 defined in the VnV Plan.
Table~\ref{tab:fr_eval} summarizes the evaluation result for each requirement.

\begin{table}[h]
\centering
\begin{tabularx}{\textwidth}{p{1.8cm}p{2.2cm}p{2.2cm}X}
\toprule
\textbf{Req.} & \textbf{Test Case(s)} & \textbf{Result} & \textbf{Evidence / Notes}\\
\midrule
R1 (BIDS input) & FR-01 & Pass & Test definition checks valid and invalid path templating, file discovery, and informative exceptions for missing files.\\
R2 (BIDS entity inference) & FR-02 & Pass & Test definition verifies subject/session/task inference from dataset structure and compares inferred values with expected entities.\\
R3 (Sensor-space analysis) & FR-03 & Pass & Test definition validates returned MNE object types and expected dimensions across filtering, epoching, and evoked-response stages.\\
R4 (Output data structures) & FR-04 & Pass & Test definition confirms source-space output type and expected source-time-course dimensions and range constraints.\\
\bottomrule
\end{tabularx}
\caption{Functional Requirements Evaluation}
\label{tab:fr_eval}
\end{table}

\section{Nonfunctional Requirements Evaluation}

\noindent Due to time constraints in the current project cycle, formal validation for usability (NFR1) and maintainability (NFR2) was not executed. Therefore, this report only includes the completed nonfunctional evaluation for portability (NFR3).

\subsection{Portability}

\begin{table}[h]
\centering
\begin{tabularx}{\textwidth}{p{2.2cm}p{2.4cm}p{2.2cm}X}
\toprule
\textbf{NFR} & \textbf{Test Case} & \textbf{Status} & \textbf{Notes}\\
\midrule
NFR3 (Portability) & NFR-03 & Pass (Validated in CI) & Cross-platform installation and execution were validated through CI runs, and the workflow completed successfully on target environments.\\
\bottomrule
\end{tabularx}
\caption{Portability Evaluation Status}
\end{table}

\section{Comparison to Existing Implementation}

Not applicable.

\section{Unit Testing}

The following matrix traces in Table~\ref{tab:unit_test} the module-level design elements (MG/MIS) to planned unit-test evidence and current status.

\begin{table}[h]
\centering
\begin{tabularx}{\textwidth}{p{2.8cm}p{2.8cm}p{2.2cm}X}
\toprule
\textbf{Module} & \textbf{Mapped Unit Tests} & \textbf{Status} & \textbf{Notes}\\
\midrule
M1 Pipeline & UT-M1-01, UT-M1-02, UT-M1-03 & Pass & Tests defined from MIS exported routines (\texttt{load\_raw}, \texttt{load\_events}, \texttt{load\_epochs}).\\
M2 TreeModel & UT-M2-01, UT-M2-02, UT-M2-03 & Pass & Covers registration, validation, and iterator semantics.\\
M3 FileTree & UT-M3-01, UT-M3-02, UT-M3-03 & Pass & Covers path expansion, error handling, and glob behavior.\\
M4 Preprocessing & UT-M4-01, UT-M4-02, UT-M4-03 & Pass & Covers load behavior and ICA workflow contracts.\\
M5 Epoch & UT-M5-01 & Pass & Verifies exported data-definition behavior.\\
M6 Covariance & UT-M6-01, UT-M6-02 & Pass & Verifies covariance generation and exception propagation.\\
\bottomrule
\end{tabularx}
\caption{Unit Test Status}
\label{tab:unit_test}
\end{table}


\section{Changes Due to Testing}

The VnV process led to the following changes in project artifacts:
\begin{itemize}
    \item Test expectations were made more explicit by adding concrete dimensional and type checks for sensor-space and source-space outputs in FR-03 and FR-04.
    \item A bug that the pipeline was unable to locate split files in the BIDS dataset was fixed.
\end{itemize}

At this stage, these changes are primarily reflected in documentation-level improvements and test-plan refinement.

\section{Automated Testing}

The VnV Plan specifies automated verification through GitHub Actions, pytest, and coverage reporting.
Automated testing has been implemented using GitHub Actions workflows, including quality checks, unit tests and system tests. The workflow files can be found in \url{https://github.com/Eelbrain/Eelbrain/tree/main/.github/workflows}


\section{Trace to Requirements}
In Table~\ref{tab:requirement_test_trace}, the requirements are traced to the system-test evidence.

\begin{table}[h]
\centering
\begin{tabularx}{0.6\textwidth}{p{3cm}X}
\toprule
    \textbf{Requirements} & \textbf{Test Cases} \\
\midrule
    R1 & FR-01 \\
    R2 & FR-02 \\
    R3 & FR-03, FR-05 \\
    R4 & FR-04, FR-05 \\
    NFR1 & NFR-01 \\
    NFR2 & NFR-02 \\
    NFR3 & NFR-03 \\
\bottomrule
\end{tabularx}
\caption{Traceability Between Requirements and System Tests}
\label{tab:requirement_test_trace}
\end{table}


\section{Trace to Modules}

In Table~\ref{tab:unit_test_module_trace}, the module-level design elements (MG/MIS) are traced to the unit-test evidence.

\begin{table}[h]
\centering
\begin{tabularx}{0.9\textwidth}{p{2.2cm}X}
\toprule
    \textbf{Module} & \textbf{Unit Test Cases} \\
\midrule
    M1 (Pipeline) & UT-M1-01, UT-M1-02, UT-M1-03 \\
    M2 (TreeModel) & UT-M2-01, UT-M2-02, UT-M2-03 \\
    M3 (FileTree) & UT-M3-01, UT-M3-02, UT-M3-03 \\
    M4 (Preprocessing) & UT-M4-01, UT-M4-02, UT-M4-03 \\
    M5 (Epoch) & UT-M5-01 \\
    M6 (Covariance) & UT-M6-01, UT-M6-02 \\
    Cross-module NFR checks & UT-NFR-01, UT-NFR-02 \\
\bottomrule
\end{tabularx}
\caption{Traceability Between Unit Tests and Modules}
\label{tab:unit_test_module_trace}
\end{table}


\section{Code Coverage Metrics}

From the statistics provided by Codecov, the code coverage for the main branch is currently at 62.56\%. See \url{https://app.codecov.io/gh/Eelbrain/Eelbrain} for the latest coverage report.


\section{Extras}

The extras set out intially (user manual and code walkthrough) are cancelled by the instructor. However, \progname{} has an actively maintained user manual at \url{https://eelbrain.readthedocs.io/en/stable/index.html}.


\bibliographystyle{plainnat}
\bibliography{../../refs/References}

\newpage{}
\section*{Appendix --- Reflection}

The information in this section will be used to evaluate the team members on the
graduate attribute of Reflection.

\input{../Reflection.text}

\end{document}